\documentclass[10pt,a4paper]{amsart}
\usepackage[margin=19mm]{geometry}
\usepackage{amsmath,amssymb,mathtools}
\usepackage[scr=rsfs,cal=euler]{mathalfa}
\usepackage{xfrac}
\usepackage[unicode,hidelinks,bookmarks=false]{hyperref}
\newcommand{\ie}{i.e.\ }
\newcommand{\cf}{cf.\ }
\newcommand{\resp}{respectively\ }
\newcommand{\Subsection}{\subsection}
\providecommand{\nobreakdash}{\nobreak\hbox{-}}
\setlength{\emergencystretch}{2em}
\multlinegap0pt
\usepackage{amssymb}
\usepackage{float}
\usepackage{tikz}
\newtheorem{lemma}{Lemma}
\def\BK{{\mathbb K}}
\title{Monomial algebras and $\mathbb{G}_a^n$-equivariant embeddings into toric varieties}
\author{Alexander Chernov}
\date{Source: arXiv:2510.22820v4 (CC BY 4.0)}
\begin{document}
\maketitle

\noindent\emph{Source and licence.} Monomial algebras and $\mathbb{G}_a^n$-equivariant embeddings into toric varieties. Version arXiv:2510.22820v4 (CC BY 4.0), distributed by arXiv under the Creative Commons Attribution 4.0 International licence, http://creativecommons.org/licenses/by/4.0/, as recorded in the arXiv licence metadata for this exact version and shown on the abstract page https://arxiv.org/abs/2510.22820v4. Full licence notice: \texttt{licenses/arxiv-2510.22820-CC-BY-4.0.txt}.\par\smallskip

\noindent\textit{Editorial scope.} Counted unit: the article's lemma lemma (source0.tex lines 332--334). The article's complete original proof of it is delivered with it (source0.tex lines 336--352, one \texttt{proof} environment of 17 lines). No mathematics is written, repaired or re-derived: the statement and the proof are the article's own lines, in the article's own notation, and no retained line is substituted or re-flowed.  No further source-local context is needed: every cross-reference of every retained line resolves inside the two delivered spans.  Nothing was omitted from any retained span, so the package carries no omission record.  No retained line contains a citation, so the wrapper carries no bibliography and no bibliography entry.  No retained line contains a figure environment, an image-inclusion command or drawing code, so no image is delivered and the package declares no asset dependency.  Editorial formatting only, in addition: an AMS \texttt{amsart} preamble that restates the article's own declarations and macros used by the retained lines, namely the environments \texttt{lemma} on separate counters, together with the standard packages those lines need.  The retained environments print in this document as Lemma 1 in order of appearance, whereas the article prints the same items with the numbering of its own full document.  The complete native source of the article as distributed in the arXiv e-print for this exact version is bundled unchanged as \texttt{sources/187840/source0.tex}.
    \begin{lemma}
        Let $B = \BK[s_1,\dotsc,s_n]/I$ be a monomial algebra with respect to a monomial ideal~$I$. Then the algebra~$\BK[\partial/\partial s_1\big|_B,\dotsc,\partial/\partial s_n\big|_B]$ is isomorphic to $B$.
    \end{lemma}

    \begin{proof}
        Denote by $\partial_i\coloneq\partial/\partial s_i|_B$. First, we prove that~$\partial_1^{k_1}\dotsc\partial_n^{k_n}= 0$ if and only if~$s_1^{k_1}\dotsc s_n^{k_n} =0$  in $B$. 

        Suppose $\partial_1^{k_1}\dotsc\partial_n^{k_n} = 0$. This immediately yields $s_1^{k_1}\dotsc s_n^{k_n} = 0$, since otherwise\begin{equation*}
            \partial_1^{k_1}\dotsc\partial_n^{k_n}(s_1^{k_1}\dotsc s_n^{k_n})= k_1!\dotsc k_n!\neq0.
        \end{equation*}
        Conversely, if $s_1^{k_1}\dotsc s_n^{k_n} = 0$ in $B$, then, since $B$ is a monomial algebra, there are no nonzero monomials $s_1^{l_1}\dotsc s_n^{l_n}$ in $B$ such that $l_i\geq k_i$. Hence if $s_1^{l_1}\dotsc s_n^{l_n}$ is a nonzero monomial in $B$, then $l_j<k_j$ for some $j$. But this yields~$\partial_j^{k_j}(s_1^{l_1}\dotsc s_n^{l_n})=0$. In particular, $\partial_1^{k_1}\dotsc\partial_n^{k_n}(s_1^{l_1}\dotsc s_n^{l_n})=0$. It means that $\partial_1^{k_1}\dotsc\partial_n^{k_n} = 0$.

        Now, note that all nonzero $\partial_1^{k_1}\dotsc\partial_n^{k_n}$ are linearly independent. Indeed, let \begin{equation*}          \sum_{\partial_1^{i_1}\dotsc\partial_n^{i_n}\neq0}\alpha_{i_1,\dotsc,i_n}\partial_1^{i_1}\dotsc\partial_n^{i_n}=0, \ \alpha_{i_1,\dotsc,i_n}\in\BK,
        \end{equation*}

        and $\alpha_{k_1,\dotsc,k_n}\neq0$ for some tuple $(k_1,\dotsc,k_n)$. Apply the derivation on the left side to the monomial $s_1^{k_1}\dotsc s_n^{k_n}$. We get \begin{equation*}
            \alpha_{k_1,\dotsc,k_n}+M = 0,
        \end{equation*} where $M$ is a linear combination of non-constant monomials and hence $\alpha_{k_1,\dotsc,k_n}=0$. This is a contradiction.

        Finally, consider the map $s_i\mapsto\partial_i$ from $B$ to $\BK[\partial_1,\dotsc,\partial_n]$. It is a well-defined homomorphism of algebras. The above reasoning shows that this map is bijective and hence it yields the required isomorphism.
    \end{proof}

\end{document}
